%! AP Statistics — Unit 1, Lesson 1.4 — Homework ANSWER KEY
\documentclass[10pt]{article}
\usepackage{apstats-article}
\usepackage{apstats-key}

%=============================================================================
\begin{document}

\pageheader{Unit 1, Lesson 1.4}{Homework --- Answer Key}
\begin{tabularx}{\linewidth}{X X X}
Name:\;\ans{Key} & Date:\; & Period:\;
\end{tabularx}

\vspace{0.1in}

% ════════════════════════════════════════════════════════════════════════════
% PART A
% ════════════════════════════════════════════════════════════════════════════
\begin{blurbbox}[Part A \quad Construct a Bar Chart]{navylight}
\small

Transportation survey ($n = 40$):

\vspace{8pt}
\renewcommand{\arraystretch}{1.8}
\begin{tabularx}{0.75\linewidth}{@{} >{\bfseries}p{0.22\linewidth}
                                        C{0.18\linewidth}
                                        C{0.26\linewidth} @{}}
\rowcolor{sky}
\textbf{Mode} & \textbf{Frequency} & \textbf{Rel.\ Frequency} \\
\midrule
Walk   & 16 & \ans{0.40 (40\%)} \\
Car    & 10 & \ans{0.25 (25\%)} \\
Bus    & 10 & \ans{0.25 (25\%)} \\
Bike   & 4  & \ans{0.10 (10\%)} \\
\midrule
\rowcolor{sky}\textbf{Total} & \ans{40} & \ans{1.00} \\
\end{tabularx}

\vspace{10pt}

\begin{teachernote}
\textbf{Correct bar chart:} 4 separated bars. X-axis: Walk, Car, Bus, Bike (any
order). Y-axis: ``Relative Frequency'' 0 to at least 0.40 (or 0 to 40\%).
Bars: Walk 0.40 (tallest), Car and Bus equal at 0.25, Bike 0.10 (shortest).
Bars do NOT touch. Title present.

Note: Car and Bus have the same height --- students should draw them identically.
\end{teachernote}

\vspace{8pt}

\begin{enumerate}[leftmargin=1.5em, itemsep=8pt]
  \item \ans{Of the 40 middle school students surveyed, walking was the most
        common mode of transportation, used by 40\% of students. Car and bus
        were equally common, each used by 25\% of students. Biking was the
        least common, at only 10\%.}

  \item \ans{Percentage \textit{not} walking $= 100\% - 40\% = 60\%$
        (or $(10 + 10 + 4)/40 = 24/40 = 0.60 = 60\%$).}

  \item \ans{Either is appropriate for 4 categories, but a bar chart is
        generally preferable here because comparing bar heights (especially
        the tie between Car and Bus at 25\% each) is clearer than comparing
        equal-sized pie slices. A pie chart could be used to emphasize that
        walking accounts for 40\% of the total (a notable fraction).}
\end{enumerate}

\end{blurbbox}

\vspace{0.14in}

% ════════════════════════════════════════════════════════════════════════════
% PART B
% ════════════════════════════════════════════════════════════════════════════
\begin{blurbbox}[Part B \quad Critique a Misleading Graph]{greenacc}
\small

\begin{enumerate}[leftmargin=1.5em, itemsep=8pt]
  \item \ans{The y-axis starts at 40 instead of 0. This means the bar for Art
        (reaching 60) looks three times taller than the bar for Music (reaching 50),
        because the visible portion of Art's bar is $60 - 40 = 20$ units and
        Music's is $50 - 40 = 10$ units. In reality, Art has only 20\% more
        students ($60/50 - 1 = 0.20$), not 100\% more as the bars imply.
        A reader incorrectly perceives that Art is dramatically more popular
        than Music.}

  \item \ans{Actual percentage: $(60 - 50)/50 \times 100\% = 10/50 \times 100\%
        = 20\%$. The headline ``Art Has 20\% More Students Than Music'' is
        numerically accurate ($60/50 = 1.20$), but it is misleading because the
        graph makes the difference appear far more dramatic than 20\%.}

  \begin{teachernote}
  Students may be surprised that the headline percentage is technically correct
  while still being paired with a misleading graph. This is a key AP Statistics
  insight: correct statistics can accompany deceptive visuals.
  \end{teachernote}

  \item \ans{Start the y-axis at 0. With a true zero baseline, both bars are
        nearly the same height (60 and 50 out of some common maximum), so the
        20\% difference is visible but not exaggerated.}

  \item \ans{Art: $60/300 = 0.200$ (20\%); Music: $50/300 \approx 0.167$ (16.7\%).
        Art has the larger share. (The remaining 190 students are in other electives.)}
\end{enumerate}

\end{blurbbox}

\vspace{0.14in}

% ════════════════════════════════════════════════════════════════════════════
% PART C
% ════════════════════════════════════════════════════════════════════════════
\begin{blurbbox}[Part C \quad AP-Style Free Response]{redacc}
\small

Car color preference ($n = 120$):

\vspace{8pt}

\begin{enumerate}[leftmargin=1.5em, itemsep=10pt]
  \item \ans{Of the 120 car buyers surveyed, white was by far the most preferred
        color (35\%), followed by black (25\%), silver (20\%), and red and
        other tied at 10\% each. White and black together account for 60\% of
        preferences. Red and other are noticeably less common, each chosen by
        only 1 in 10 buyers.}

  \begin{teachernote}
  Full credit: $n$ stated; variable in context; most and least common categories
  named with relative frequencies; at least one additional observation (e.g.,
  combined frequencies, comparison).
  \end{teachernote}

  \item \ans{Sector angles:
        White $= 0.350 \times 360 = 126^\circ$;
        Black $= 0.250 \times 360 = 90^\circ$;
        Silver $= 0.200 \times 360 = 72^\circ$;
        Red $= 0.100 \times 360 = 36^\circ$;
        Other $= 0.100 \times 360 = 36^\circ$;
        Total $= 360^\circ$.\\[4pt]
        Black and Silver together: $90 + 72 = 162^\circ$, which is almost
        exactly half the circle ($180^\circ$). Together they represent
        $25\% + 20\% = 45\%$ of preferences.}

  \begin{teachernote}
  Also accept ``White and Black'' ($126 + 90 = 216^\circ$, which is 60\% of
  the circle, not half). The closest pair to exactly half the circle is Silver
  + Red + Other ($72 + 36 + 36 = 144^\circ$, 40\%) or Black + Silver ($162^\circ$,
  45\%). Accept any reasonable reading of ``almost exactly half.''
  \end{teachernote}

  \item \ans{A bar chart is better for a dealership poster because it makes it
        easy to compare the relative popularity of each color side by side.
        With 5 categories, bar heights are simpler to compare than pie sector
        angles. A pie chart would work but makes the equal Red and Other slices
        harder to distinguish.}
\end{enumerate}

\end{blurbbox}

\end{document}
